\section{Conclusion and Future Work}
\label{sec:conclusion}

This paper introduced AvocadoLeaf, a 1,861-image dataset of five field-collected avocado leaf disease categories and healthy leaves, and used it to benchmark 15 lightweight and efficient architectures for edge-deployable disease classification. ConvNeXt‑V2 offered the best accuracy-efficiency trade-off, reaching 92.14\% accuracy and 0.900 macro F1 with only 3.4M parameters and 2.30 ms per-image inference on a single GPU, though its accuracy is statistically tied with five other architectures (Section~\ref{subsec:arch_compare}, \ref{app:bootstrap}). Metric-based few-shot learning with a frozen ConvNeXt‑V2 backbone more than doubled the accuracy of conventional fine-tuning at every shot level tested (e.g. 57.4\% vs. 28.6\% at 20 shots), and uncertainty estimation showed that the model reliably separates in-distribution avocado leaves from both unrelated images (far OOD AUC 0.994) and leaves of other crop species (near OOD AUC 0.985), while a full-test-set LIME deletion-metric evaluation confirmed that its predictions localise to disease-relevant leaf regions rather than background artefacts (mean confidence drop 0.649 when the attributed region is hidden, 95\% CI [0.618, 0.678], $n=280$). Building on this, we combined the CNN's prediction with LIME's visual attribution and a small, free, locally-hosted vision-language model to generate grounded, natural-language explanations of individual predictions, while keeping the CNN as the sole decision-maker (Section~\ref{subsec:vlm_explain}). Together, these results indicate that an accurate, compact, and interpretable classifier can be built for avocado disease diagnosis under realistic, field-collected conditions.

We additionally examined several deployment-facing questions raised by these results. First, a naive-deployment failure-mode analysis (Section~\ref{subsec:cross_dataset}) showed that ConvNeXt‑V2's advantage over the lighter candidates is not limited to in-distribution accuracy. Without an uncertainty gate, MobileNet‑V2 confidently (${>}0.9$) misclassifies 22--30\% of domain-shifted, non-avocado images as a specific disease, versus 0--1.5\% for ConvNeXt‑V2, making the entropy-based uncertainty gate (Section~\ref{subsec:uncertainty}) a necessary companion to any lighter model chosen for deployment rather than an optional add-on. Second, and more directly, we obtained a genuine independent, same-task cross-domain evaluation (Section~\ref{subsec:real_cross_domain}) using K-Kotagiri \citeyearpar{kotagiri2024avocado}, a separately collected avocado leaf dataset from India: trained only on AvocadoLeaf and evaluated zero-shot at the healthy-vs-diseased level (the only granularity Kotagiri's labels support), accuracy falls from 98.57\% in-domain to 60.00\%, and a few-shot adaptation fallback using as few as 20 labelled target-domain images per class recovers only part of this gap (up to 68.78\%). This is the single most important caveat on this paper's generalisation claims: strong in-distribution accuracy does not, on its own, indicate that BoVien is ready to deploy at a new site without either site-specific labelled data or further adaptation. Third, we deployed BoVien live across four platforms (Section~\ref{subsec:realtime_demo}), running natively on-device on Android (ONNX Runtime, CPU only) and iOS (native Core ML) as well as via ONNX Runtime on Windows and macOS laptop compute, obtaining real per-platform confidence, FPS, and latency numbers from a live camera rather than only a batched GPU benchmark. Fourth, beyond identifying gaps, we resolved one of them directly: post-hoc temperature scaling reduces test-set ECE from 0.0311 to 0.0193 (Section~\ref{subsec:uncertainty}) at no cost to accuracy. We also characterised, rather than merely flagged, where BoVien's remaining errors concentrate: a confusion-matrix-driven error analysis (Section~\ref{subsec:error_analysis}) shows they cluster in the three lowest-support classes and in three visually-motivated confusion pairs (\textit{gisat}$\leftrightarrow$\textit{moctrang}, \textit{repsap}/\textit{moctrang}$\to$\textit{chaymepla}), consistent with the LIME faithfulness results (Section~\ref{subsec:interpretability}).

At the same time, Section~\ref{sec:discussion} identified concrete gaps that remain open. A dedicated embedded board such as a Raspberry Pi remains untested, and the on-device numbers above exclude the background-removal preprocessing step, so true end-to-end field latency would be somewhat higher, particularly on Android. The real cross-domain evaluation above is same-task but binary-only, so disease-level (6-class) cross-domain generalisation is still untested. The ablation study (Section~\ref{subsec:ablation}) also surfaced a documentation/protocol discrepancy in the shared augmentation policy across the 15-architecture benchmark. We also tested, and found unsuccessful, a harder use of the same language model beyond explanation generation, using its explanation-derived judgement as a standalone error-detection signal, which underperformed the CNN's own existing uncertainty estimation (Section~\ref{sec:discussion}, \ref{app:llm}). Future work will therefore prioritise (1) Raspberry-Pi/embedded-board latency measurement that includes the background-removal step, (2) collecting an independent, same-task, \emph{6-class} AvocadoLeaf dataset to extend the binary cross-domain result above to disease-level generalisation, and (3) a fully protocol-matched re-run of the 15-architecture benchmark under a uniform augmentation policy to close the discrepancy noted above. Beyond these, we plan to expand the dataset to more growing regions, investigate why RelationNet failed to learn a useful metric in the few-shot setting in contrast to ProtoNet, validate the VLM-generated explanations against expert (agronomist) judgement rather than only the automated correctness proxy used here, and conduct a more extensive, systematic exploratory data analysis to guide future data collection toward the classes and conditions the model currently finds hardest.
